\section{Related Work}
\label{sec:relwork}

\paragraph{Proof rules for almost-sure termination.}
Classical bounded-variant arguments establish AST by requiring a uniformly positive probability of descent within a bounded natural-valued range~\cite{DBLP:series/mcs/McIverM05}. Supermartingale-based rules extend the scope of this approach. McIver et al.~\cite{10.1145/3158121} formulate an AST rule directly over source programs, including demonic choice, using one real-valued function with both supermartingale and progress obligations. At pCFG level, Majumdar and Sathiyanarayana~\cite{DBLP:journals/pacmpl/MajumdarS25} establish sound and relatively complete proof rules and also prove relative completeness of the earlier rule of McIver et al.

The two-certificate rule of Majumdar and Sathiyanarayana is our starting point: it already separates the supermartingale from the natural-valued variant. We transfer this structure to complete loop iterations in distribution-transformer semantics and relate the resulting certificates to pCFG certificates in both directions. Relative to their rule, the contribution is the program-level formulation, its metatheory, and the reduction in location-dependent proof obligations. Relative to McIver et al., the distinction is the use of two separately chosen functions, coupled through sublevel boundedness, instead of requiring one function to serve both roles. This permits a variant to reset on rejection without itself satisfying a supermartingale inequality. Our FDR comparison makes the change in certificate representation explicit; FLDR and Han--Hoshi illustrate applications of the source-level rule. We address qualitative AST, which permits infinite expected runtime.

\paragraph{Operational and denotational models.}
pCFGs and Markov decision processes are the standard operational substrate for termination and reachability analysis of probabilistic and randomized systems~\cite{DBLP:books/wi/Puterman94,DBLP:books/sp/BaierK08}. Tools such as Storm~\cite{DBLP:journals/sttt/HenselJKQV22} and refined analyses of visiting times and stationary behavior~\cite{DBLP:conf/tacas/MertensKQW24,DBLP:journals/jcss/JungesK0W21} operate at this level. Viewing probabilistic systems as \emph{distribution transformers} rather than as single-state transition systems has recently been fruitful for both Markov chains and MDPs: Aghamov et al.~\cite{DBLP:journals/corr/abs-2406-15087,DBLP:conf/birthday/AghamovBKNOPV25} model-check Markov chains through their action on distributions, and Akshay et al.~\cite{DBLP:conf/cav/AkshayCMZ23,DBLP:conf/ijcai/0001CMZ24} synthesize invariants and certified policies for MDPs under distributional reach-avoidance objectives; distribution-based objectives for MDPs are also studied in~\cite{DBLP:conf/lics/AkshayGV18}. Our rule applies the same distribution-transformer perspective one level higher, directly to program syntax rather than to an underlying automaton, which avoids the location-indexed bookkeeping that a translation to pCFGs or MDPs otherwise imposes.

\paragraph{Denotational semantics and program logics for probabilistic programs.}
The denotational, post-distribution-transformer semantics we build on originates with Kozen~\cite{DBLP:conf/focs/Kozen79,DBLP:conf/stoc/Kozen83}. It also underlies the calculi for weakest preconditions and weakest pre-expectations of McIver and Morgan~\cite{DBLP:series/mcs/McIverM05} and Kaminski~\cite{DBLP:phd/dnb/Kaminski19}, in the tradition of Dijkstra's predicate-transformer semantics~\cite{DBLP:journals/cacm/Dijkstra75,DBLP:books/daglib/0067387}. Hoare-style program logics for probabilistic programs were introduced by den Hartog and de Vink~\cite{DBLP:journals/ijfcs/HartogV02} and extended by the Ellora framework~\cite{DBLP:conf/esop/BartheEGGHS18}; more recent calculi include quantitative strongest-postcondition reasoning~\cite{DBLP:journals/pacmpl/ZhangK22}, its hyperproperty generalization~\cite{DBLP:journals/corr/abs-2404-05097}, and modal separation logic for conditional probability~\cite{DBLP:journals/pacmpl/Li0H23}. These frameworks provide broader forms of reasoning about probabilistic behavior. Our focus is a specific AST rule that can be used alongside distributional partial-correctness arguments.

\paragraph{Verifying exact sampling algorithms.}
Exact discrete samplers built from fair coin flips, following the classical constructions of von Neumann~\cite{vonNeumann1951} and Lumbroso's Fast Dice Roller (FDR)~\cite{DBLP:journals/corr/abs-1304-1916}, have received growing verification attention. Saad et al.\ introduce the Fast Loaded Dice Roller (FLDR) as a near-optimal sampler for finite discrete distributions~\cite{DBLP:conf/aistats/SaadFRM20}. Han and Hoshi's interval algorithm~\cite{HanHoshi1997} supplies the basis for our third sampler case study. Bagnall, Stewart, and Banerjee's Zar compiler~\cite{DBLP:journals/pacmpl/Bagnall0023,10.1145/3591220} formally verifies samplers compiled from probabilistic programs with loops and conditioning in Coq, targeting the random-bit model; generating-function techniques have likewise been used to reason about probabilistic programs and about samplers such as the fast dice roller~\cite{DBLP:conf/lopstr/KlinkenbergBKKM20,DBLP:conf/cav/ChenKKW22,haase2026generatingfunctionsmeetoccupation}. Exact and guaranteed-bound inference engines such as PSI~\cite{DBLP:conf/cav/GehrMV16,DBLP:conf/pldi/GehrSV20}, the approach of Beutner et al.~\cite{DBLP:conf/pldi/BeutnerOZ22}, and scalable exact inference for discrete programs~\cite{DBLP:journals/pacmpl/HoltzenBM20} address the complementary problem of computing or bounding the output distribution.

Closest to our work is the distributional-invariant framework of Zilken et al.~\cite{zilken2026verifyingsamplingalgorithmsdistributional}. It already establishes partial and total correctness of FDR and FLDR, with AST arguments using existing rules in its appendices. We build on those invariants and verification examples to give a general program-level AST rule and an explicit comparison of proof representations. The Han--Hoshi case study adds a different sampling mechanism, using interval boundaries as the termination certificate. Its role here is to demonstrate applicability of the AST rule; the present manuscript does not develop a separate distributional output-correctness proof for Han--Hoshi.

\paragraph{Scope of the synthesis contribution.}
The prototype in \Cref{sec:certificate-synthesis} instantiates our AST
obligations with template-based counterexample-guided search. Its contribution
here is the shared treatment of both certificate roles over whole iterations,
together with explicit assumptions and reproducible examples. We do not claim
CEGIS, affine templates, or the separation of the two roles as new techniques.
The random walk distinguishes qualitative AST from proofs requiring finite
expected runtime. The FDR result, by contrast, already supplies a ranking
supermartingale. The two FLDR experiments distinguish exact-instance synthesis
from synthesis conditional on a structural abstraction, and show that a
separation within a restricted feature basis need not survive a richer basis.
These examples support the applicability of our program-level formulation;
they do not establish a general advantage over existing certificate-synthesis
methods or over the single-function AST rule of McIver et al.~\cite{10.1145/3158121}.

\paragraph{Nondeterminism.}
Several lines of work verify probabilistic programs or Markov models in the presence of demonic or angelic nondeterminism~\cite{10.1145/3158121,DBLP:journals/pacmpl/BatzBKW24,Wang2018ADS,Zilberstein2025,10.1145/3743131,Zilberstein_2025}. Demonic AST requires termination with probability one under every admissible scheduler; both McIver et al.~\cite{10.1145/3158121} and the pCFG framework of Majumdar and Sathiyanarayana~\cite{DBLP:journals/pacmpl/MajumdarS25} support this interpretation. Our rule instead uses a single post-distribution for each input and currently excludes nondeterministic choice. Extending it requires a strategy-indexed or set-valued semantics and corresponding invariant and progress obligations. Strategy synthesis, as studied by Batz et al.~\cite{DBLP:journals/pacmpl/BatzBKW24}, offers a complementary route for proving AST under a selected strategy. We discuss the precise obstacles and possible extensions in \Cref{sec:conclusion}.
